% s11_final (abridged version). The resolved conifold, the evidence, and what the results are conditional on.
\section{The resolved conifold, the evidence, and what remains conditional}
\label{sec:fin}
\noindent\genusstatus{numerically supported} The finite-order checks recorded in this section are computations (exact arithmetic unless stated); they support the statements at the orders and genera given, and are not proofs. The code and outputs are in \texttt{code/} of this project.


\emph{The resolved conifold.} \Lg{Section~11} records the counterpart of the preceding results for the resolved conifold times $\bbC^2$, the total space of $\cO(-1)^{\oplus 2}\oplus\cO^{\oplus 2}$ over $\bbP^1$, on its fibrewise torus. That section of \textup{[L]} is not in final form, and what follows is conditional: its central statement, the rank-$2$ leaf structure \Lg{Theorem~11.1}, is stated in \textup{[L]} without proof, and the regularity statement \Lg{Corollary~11.4} depends on it.

With $\epsilon_1=\xc$, $\epsilon_2=1$ and $\cc\coloneqq1+\xc$, $\xc\cc\neq0$, one has $E=\cc$ and $\bw=(\cc,-\xc,-1)$; the coordinate is $\tcn\coloneqq\log\cQ$ and the conifold point is $\tcn=0$. There are two colours, and the frame-B factors and the decoration sums are defined in \Lg{Section~11.1} on the pattern of \cref{sec:set}. The analogue of \cref{thm:lea-main} is the rank-$2$ leaf structure \Lg{Theorem~11.1}. No proof of it is given in \textup{[L]} or here. \textup{[L]} only names the ingredients of a proof, by a route different from that of \cref{sec:lea}: an Euler integral and a closed form of the singular row through Gauss's hypergeometric function, the quantum differential equation at each idempotent, a flow equation for the edge at the singular idempotent, and one analytic input, Stirling's series with remainder \cite[5.11.1, Section~5.11(ii)]{DLMF}. The vanishing $\Mirr(h)$ and the value of the vertex do not depend on this: they follow from \cref{cor:van-scale} and \cref{thm:val-tau} at $E=\cc$ and $\tau=-\xc/\cc$. Given \Lg{Theorem~11.1}, the argument of \cref{thm:reg-main,cor:reg-sv} gives the following.

\begin{cor}[Regularity for the conifold, conditional; {\Lg{Corollary~11.4}}]\genusstatus{AI-proposed, unchecked}
\label{cor:con-reg}
Assume the rank-$2$ leaf structure \Lg{Theorem~11.1}. Then for every $g\geq 2$ and every $\xc$ with $\xc\cc\neq 0$: every coloured stable graph of genus $g$ with at least one edge is regular at $\tcn=0$ after summing over decorations, and the polar part of the single-vertex graph is $\beta_{2g-2}(\xc,1)\,\tcn^{2-2g}$.
\end{cor}

In \textup{[L]} the corollary is worded without this hypothesis, and its proof starts from \Lg{Theorem~11.1}; so it is not proved as long as that theorem is not. The identification of the graph sum with the generating series of the conifold is \Lg{Proposition~11.7}. Its proof in \textup{[L]} follows \cref{sec:gen,sec:fra} in rank $2$; several steps are written by reference to the rank-$3$ proofs and to a formal Laplace calculus in a separate notation, and \textup{[L]} leaves open whether they will be written out in rank $2$. It uses \cite{Tel12,GP99,CG07,PPZ15}, Mumford's formula and \cite[Theorem~4]{FP00}, the Aspinwall--Morrison formula as in \cite[p.~3]{FP00}, and Stirling's series \cite[5.11.1, 5.12.1]{DLMF}; it does not use \cite[Theorem~31]{CCIT15}, and \cref{cor:con-reg} does not use it. The ratio $(-3)^{g-1}$ between the two polar parts comes from the prefactors of $\Delta$ at the singular colour, $3\epsilon_1\epsilon_2E^3$ against $-\xc\cc^3$; the vertex value is the same \Lg{Remark~10.16}.

\emph{Finite-order evidence.} \Lg{Appendix~A} collects finite-order checks. Nothing in it is a proof step, and no proof uses any of it. The arithmetic is exact unless stated otherwise there; each check compares two quantities computed in different ways, and many are accompanied by controls, that is, deliberately altered versions of the computation which are observed to fail. The checks concern the following links of the argument: the graph sum of \cref{thm:gen-main} against virtual localisation in genus $2$, $3$, $4$ at low degree; the closed form of the rows against the Wick expansion of $\Phi$; the specialisation and the frame change at genus $2$, $3$, $4$ at sample points; the gap itself, in that the totals at genus $2$, $3$, $4$ at $\sfF=\Fcf$, as exact Laurent series in $\xi$ at five weight pairs, have no polar coefficient other than that of \cref{cor:int-gap}, and per colouring all graphs with an edge are regular at genus $2$ and $3$ at two weight pairs; the leaf structure at finite orders; \cref{thm:van-main}, also directly on the Hodge side; and \cref{thm:val-tau} for $h\leq8$ from Mari\~no--Vafa data and, independently of the Mari\~no--Vafa formula, for $h\leq5$. \Lg{Section~A.7} states the limits of this evidence. Nothing was checked to all orders by computation. Most items are at $g\leq3$; genus $4$ is covered at one or two weight pairs and points; at $g\geq5$ no full total has been computed. Some controls do not discriminate: sign controls of rank one at odd genus, and tests of first order in $z$ between $\Rm(z)$ and $\Rm^{-1}(-z)$. The identity $\Omega=\Rm.\TFT$ at the level of classes was paired only with $\psi$ classes and powers of $\lambda_1$ at low degree. At genus $2$, $3$, $4$ the computed $\sF_{k,h}$ have coefficients in $\bbQ[X^{\pm1}]$; this is evidence only for the stronger form of \cref{sec:int-notclaimed}.

\emph{What the results are conditional on.} For the main statements, \cref{thm:int-frame,thm:int-regularity,cor:int-gap,cor:int-gapkh}, no hypothesis is left unproved in \textup{[L]}: the hypothesis $\Mirr(h)$ of \cref{thm:reg-main,cor:reg-sv} is established by \cref{lem:reg-A}\textup{(b)}. These results are conditional in the following sense.
\begin{itemize}
\item \emph{Quoted theorems.} The published results of \cref{sec:int-inputs} are used as stated and are not reproved. \Cref{thm:int-regularity} rests on the Mari\~no--Vafa formula \cite{LLZ03} and textbook facts alone. \Cref{thm:int-frame} rests on the mirror theorem \cite[Theorem~31]{CCIT15}, on \cite{Tel12} with the stability results behind it, on \cite{PPZ15,GP99,CG07,FP00,Mum83,Pan98}, and on \cite[Eq.~15.6.7]{DLMF} with Gauss's values. \Cref{cor:int-gap,cor:int-gapkh} rest on all of these and on Barnes's second lemma \cite{Bar10}. Beyond that list, \cite[Lemma~4.2]{BFGW21} is used only for statements involving $A$ (\cref{sec:fra}), and the inputs named above for the resolved conifold are used only there.
\item \emph{Checks which \textup{[L]} marks as outstanding.} The long version is a draft, and it flags the following. The numbers quoted from \cite{CCIT15} and \cite{Tel12} are those of the arXiv versions and have not been compared with the journal versions. The formulas of the action of \cite{PPZ15} recalled in \cref{sec:set-hodge} are written in the form in which the proofs use them, and have not been compared with the printed definitions, including the sign in the symplectic condition. The form of Barnes's second lemma used in \cref{lem:fra-tcf} was confirmed only through secondary statements, not from a printed source. In \cref{sec:van}, the pointers to the textbooks for the principal specialisation of Schur functions and for the hook lengths of a row are written from memory, and no theorem numbers are fixed for the Frobenius formula, the Murnaghan--Nakayama rule, Frobenius' transposition identity, or the string and dilaton equations with pulled-back classes. The axioms of the virtual class are used in standard form with no reference given. The references \cite{HK10,CI18,ZZ08} of \cref{sec:int-problem} are quoted after \cite{BS} and were not consulted; of \cite{LP18} and \cite{BFGW21} only the formulas compared and Lemma~4.2, respectively, were read. Until these checks are made, the corresponding steps rest on the forms of the quoted results as written in \textup{[L]}.
\item \emph{Readings adopted as definitions.} The conifold gap is a statement about the substitution $\Cf$ of \cref{def:fra-cf}, a substitution on the unique polynomial in $F$ and not an analytic continuation of the series; this is the reading of \cite{BS} adopted here. The branch of $\log y$ for $\tcf$ is that of \cref{lem:fra-tcf}, and in \cref{cor:int-gapkh} the period $\Pi_{\gamma_{\mathrm{cf}}}$ is read through its expansion.
\item \emph{Scope.} All statements are on the fibrewise torus with $\epsilon_1\epsilon_2E\neq0$. The points listed in \cref{sec:int-notclaimed} are not claimed.
\item \emph{The resolved conifold.} \Cref{cor:con-reg} is conditional on the rank-$2$ leaf structure \Lg{Theorem~11.1}, which is stated without proof in \textup{[L]}, and the proof of \Lg{Proposition~11.7} is written there by reference to the rank-$3$ proofs. Neither is used in the main statements.
\end{itemize}
